\documentclass{beamer}

\mode<presentation>
\title{SoC1\&2: Formal Equivalence with Conformal LEC}
\author[AS]{Andreas Siggelkow \newline \texttt{andreas.siggelkow@rwu.de}}
\date[2026]{\today}
\subtitle{RWU}
\institute{Elektrotechnik und Informatik}

%%%%%%%%%%%%%%%%%%%%%%%%%%%%%%%%
% Logo
%%%%%%%%%%%%%%%%%%%%%%%%%%%%%%%%
\pgfdeclareimage[height=1cm]{hslogo}{./hs}
\logo{\pgfuseimage{hslogo}}
%%%%%%%%%%%%%%%%%%%%%%%%%%%%%%%%
% Macros
%%%%%%%%%%%%%%%%%%%%%%%%%%%%%%%%
\usepackage[T1]{fontenc}
\usepackage{lmodern}
\usepackage{textcomp}
\definecolor{links}{HTML}{2A1B81}
\hypersetup{colorlinks,linkcolor=,urlcolor=links}
\usepackage{graphicx}
\usepackage{ragged2e, tabularx}
\newcolumntype{Y}{>{\RaggedRight\arraybackslash}X}
\usetheme{Singapore}
\usecolortheme{seahorse}
\usefonttheme[onlylarge]{structuresmallcapsserif}
\usefonttheme[onlysmall]{structurebold}
\setbeamertemplate{footline}[frame number]
\usepackage{listings}
\newcommand{\astitle}{nix}
\newcommand{\assubtitle}{nix}
\newcommand{\asfolie}[3]{\begin{frame}\frametitle{#1}\begin{block}{#2}{#3}\end{block}\end{frame}}

% Commands and tool output are set with listings: dashes, $, _, \ and quotes
% appear exactly as they have to be typed.
\definecolor{lstbg}{HTML}{F2F2F7}
\definecolor{lstcomment}{HTML}{5A5A8A}
\lstset{
  basicstyle=\ttfamily\scriptsize,
  columns=fullflexible,
  keepspaces=true,
  upquote=true,
  breaklines=true,
  breakatwhitespace=false,
  backgroundcolor=\color{lstbg},
  frame=none,
  xleftmargin=2pt,
  xrightmargin=2pt,
  aboveskip=4pt,
  belowskip=4pt,
  showstringspaces=false
}
\lstdefinestyle{sh}{
  morecomment=[l]{\#},
  commentstyle=\color{lstcomment}\itshape
}
\lstdefinestyle{out}{basicstyle=\ttfamily\tiny}
\newcommand{\cmd}[1]{\texttt{#1}}




\begin{document}
%%%%%%%%%%%%%%%%%%%%%%%%%%%%%%%%
% First pages
%%%%%%%%%%%%%%%%%%%%%%%%%%%%%%%%
\setcounter{tocdepth}{1}

\begin{frame}
  \titlepage
\end{frame}

\begin{frame}
  \frametitle{Outline}
  \tableofcontents
\end{frame}

%%%%%%%%%%%%%%%%%%%%%%%%%%%%%%%%
% Overview
%%%%%%%%%%%%%%%%%%%%%%%%%%%%%%%%
\section{Overview}
\renewcommand{\astitle}{Overview}

\renewcommand{\assubtitle}{What Is an Equivalence Check?}
\begin{frame}
\frametitle{\astitle}
\begin{block}{\assubtitle}
  \begin{itemize}
    \item Two descriptions of the same design: \emph{golden} (reference)
      and \emph{revised} (to be checked)
    \item Conformal LEC \emph{proves} that both compute the same function
      -- for every input value, without a testbench and without
      simulation time
    \item Simulation shows: these 62 tests pass. LEC shows: no input
      sequence at all can tell RTL and netlist apart
    \item It does not show that the RTL is right: a bug in the RTL is also
      in the netlist, and LEC calls them equivalent
  \end{itemize}
\end{block}
\end{frame}

\renewcommand{\assubtitle}{Equivalence vs. Property Checking}
\begin{frame}
\frametitle{\astitle}
\begin{block}{\assubtitle}
  \footnotesize
  \begin{tabularx}{\linewidth}{@{}lYY@{}}
    & \textbf{equivalence (LEC)} & \textbf{property (Jasper, SymbiYosys)}\\\hline
    question & design A $\equiv$ design B? & does the design meet assertion P?\\
    needs    & two designs              & one design + SVA properties\\
    finds    & tool or flow errors      & design errors\\
    tool     & Conformal LEC            & Jasper (no licence), SymbiYosys\\
    effort   & low: scripts only        & high: write the properties\\
  \end{tabularx}
  \vspace{4pt}
  \begin{itemize}
    \item LEC is part of every sign-off: it closes the gap between the
      simulated RTL and the netlist that goes to the fab
  \end{itemize}
\end{block}
\end{frame}

\renewcommand{\assubtitle}{Where in the Flow?}
\begin{frame}
\frametitle{\astitle}
\begin{block}{\assubtitle}
  \footnotesize
  \begin{tabularx}{\linewidth}{@{}lp{0.36\linewidth}Y@{}}
    \textbf{check} & \textbf{golden $\rightarrow$ revised} & \textbf{what can go wrong}\\\hline
    \cmd{rtl2syn} & RTL $\rightarrow$ Genus netlist (\cmd{\_syn.v})
                  & optimization, constant propagation, scan insertion,
                    Verilog read differently than by the simulator\\
    \cmd{syn2pnr} & Genus netlist $\rightarrow$ Innovus netlist
                    (\cmd{\_routed.v})
                  & CTS buffers, resizing, hold fixing, scan reorder\\
  \end{tabularx}
  \vspace{4pt}
  \begin{itemize}
    \item Both together: RTL $\equiv$ layout netlist
    \item Run \cmd{rtl2syn} after every synthesis, \cmd{syn2pnr} after
      every P\&R run
    \item Faster than a gate-level regression, and complete
  \end{itemize}
\end{block}
\end{frame}

%%%%%%%%%%%%%%%%%%%%%%%%%%%%%%%%
% How LEC works
%%%%%%%%%%%%%%%%%%%%%%%%%%%%%%%%
\section{How LEC Works}
\renewcommand{\astitle}{How LEC Works}

\renewcommand{\assubtitle}{Key Points}
\begin{frame}
\frametitle{\astitle}
\begin{block}{\assubtitle}
  \begin{itemize}
    \item LEC does not unroll time: it cuts the design at the
      \emph{key points}
    \begin{itemize}
      \item primary inputs and outputs
      \item flip-flops (DFF) and latches (DLAT)
      \item black boxes (BBOX): here the SPRAM macros
      \item cut points (added on purpose, see DFT)
    \end{itemize}
    \item Between the key points there is only combinational logic
    \item Two steps:
    \begin{enumerate}
      \item \textbf{mapping}: which flip-flop in the netlist is which
        flip-flop in the RTL (by name, then by function)
      \item \textbf{compare}: is the logic cone in front of each mapped
        key point the same Boolean function?
    \end{enumerate}
    \item Needs the same states in both designs (no retiming here)
  \end{itemize}
\end{block}
\end{frame}

\renewcommand{\assubtitle}{Mapping: Unmapped Points}
\begin{frame}
\frametitle{\astitle}
\begin{block}{\assubtitle}
  \footnotesize
  \begin{tabularx}{\linewidth}{@{}lY@{}}
    \textbf{class} & \textbf{meaning}\\\hline
    unreachable & no output depends on it -- Genus removed it, the RTL
                  still has it: harmless\\
    extra       & only in one design, e.g. scan logic, CTS buffers
                  are not key points; DFT lockup latches\\
    not-mapped  & in both designs, but LEC found no partner:
                  \textbf{look at it}\\
  \end{tabularx}
  \vspace{4pt}
  \begin{itemize}
    \item \cmd{lec\_<mode>\_unmapped.rpt}
    \item Not-mapped points become non-equivalent points later
    \item Typical cause: a flip-flop renamed beyond the LEC name rules,
      or two flip-flops merged by Genus (\cmd{set\_analyze\_option
      -auto} handles most cases)
  \end{itemize}
\end{block}
\end{frame}

\renewcommand{\assubtitle}{Compare Result}
\begin{frame}
\frametitle{\astitle}
\begin{block}{\assubtitle}
  \footnotesize
  \begin{tabularx}{\linewidth}{@{}lY@{}}
    \textbf{result} & \textbf{meaning}\\\hline
    equivalent     & proven equal\\
    non-equivalent & LEC found an input combination with different
                     values: a counterexample exists\\
    abort          & the solver gave up (large arithmetic cones);
                     neither proven nor refuted\\
  \end{tabularx}
  \vspace{4pt}
  \begin{itemize}
    \item Sign-off: \textbf{0 non-equivalent, 0 abort}
    \item \cmd{analyze\_abort -compare} retries aborts with more effort
    \item Black-box pins (SPRAM) are compared like outputs: address,
      data, write enable of the macro must be the same
  \end{itemize}
\end{block}
\end{frame}

%%%%%%%%%%%%%%%%%%%%%%%%%%%%%%%%
% DFT
%%%%%%%%%%%%%%%%%%%%%%%%%%%%%%%%
\section{Scan Chain}
\renewcommand{\astitle}{Scan Chain}

\renewcommand{\assubtitle}{The Netlist Has More Than the RTL}
\begin{frame}
\frametitle{\astitle}
\begin{block}{\assubtitle}
  \begin{itemize}
    \item Genus replaced every core flip-flop by a scan flip-flop:
      $D' = SE\;?\;SI : D$
    \item SE is \cmd{core/as\_jtag/sc01\_shift\_o} (TAP, SCANTEST +
      Shift-DR); in the RTL nothing uses it
    \item The chain end goes to \cmd{core/as\_jtag/sc01\_tdo\_i}; in the
      RTL (\cmd{XFAB\_DFT}) this pin is a constant 0
    \item Without constraints: every scan flip-flop is non-equivalent,
      and TDO in SCANTEST as well
  \end{itemize}
\end{block}
\end{frame}

\renewcommand{\assubtitle}{Constraints in \texttt{04\_rv64\_lec.tcl}}
\begin{frame}[fragile]
\frametitle{\astitle}
\begin{block}{\assubtitle}
\begin{lstlisting}[style=sh]
# shift enable: cut, becomes a primary input, held at 0
add_primary_input core/as_jtag/sc01_shift_o -pin -cut -both
add_pin_constraints 0 core/as_jtag/sc01_shift_o -both
# chain end: cut in both designs, a free input on both sides
add_primary_input core/as_jtag/sc01_tdo_i   -pin -cut -both
\end{lstlisting}
  \begin{itemize}
    \item The internal TAP pins are no ports: a \emph{cut} makes them a
      key point LEC may constrain
    \item With SE = 0 the scan multiplexers are transparent: the
      functional mode is compared, completely
    \item The scan mode itself is checked by the gate-level simulation
      (test \cmd{scan}) and later by ATPG
    \item Same constraints in \cmd{syn2pnr}: a scan reorder in Innovus
      would otherwise change the chain
  \end{itemize}
\end{block}
\end{frame}

%%%%%%%%%%%%%%%%%%%%%%%%%%%%%%%%
% Scripts
%%%%%%%%%%%%%%%%%%%%%%%%%%%%%%%%
\section{Scripts}
\renewcommand{\astitle}{Scripts}

\renewcommand{\assubtitle}{Files}
\begin{frame}
\frametitle{\astitle}
\begin{block}{\assubtitle}
  \footnotesize
  \begin{tabularx}{\linewidth}{@{}lY@{}}
    \textbf{file} & \textbf{role}\\\hline
    \cmd{scripts/lec/lec.sh}        & wrapper: inputs, run, PASS/FAIL\\
    \cmd{scripts/04\_rv64\_lec.tcl} & Conformal Tcl mode: read, constrain,
                                      map, compare, reports\\
    \cmd{scripts/rv64\_rtl\_setup.tcl} & RTL file list + defines, shared
                                      with Genus (\cmd{00\_rv64\_elab.tcl})\\
    \cmd{build\_genus/cc\_chip\_top\_rtl2syn\_genus.do} & dofile from
                                      Genus \cmd{write\_do\_lec} (option
                                      \cmd{-g})\\
    \cmd{build\_lec/}               & logs, reports (\cmd{.gitignore})\\
  \end{tabularx}
  \vspace{4pt}
  \begin{itemize}
    \item The golden RTL must be read \emph{exactly} as Genus read it:
      same files, same defines (\cmd{XFAB\_SPRAM}, \cmd{XFAB\_PADS},
      \cmd{XFAB\_CLKCELLS}, \cmd{XFAB\_DFT}, cell names) -- hence one
      shared file
  \end{itemize}
\end{block}
\end{frame}

\renewcommand{\assubtitle}{\texttt{04\_rv64\_lec.tcl} (short)}
\begin{frame}[fragile]
\frametitle{\astitle}
\begin{block}{\assubtitle}
\begin{lstlisting}[style=sh]
source $::env(XFAB_PATHS)
source $SCRIPT_DIR/rv64_rtl_setup.tcl        ;# RTL_FILES, RTL_DEFINES
read_library -both -liberty $LIBS_TYP
read_design -golden  -sv09 -define ... $RTL_FILES   ;# rtl2syn
read_design -revised -verilog $SYN_NETLIST
set_root_module cc_chip_top -golden
set_root_module cc_chip_top -revised
add_primary_input core/as_jtag/sc01_shift_o -pin -cut -both
add_pin_constraints 0 core/as_jtag/sc01_shift_o -both
add_primary_input core/as_jtag/sc01_tdo_i   -pin -cut -both
set_flatten_model -seq_constant -seq_constant_x_to 0
set_analyze_option -auto
set_system_mode lec                          ;# mapping
add_compared_points -all
compare
analyze_abort -compare
report_verification > lec_rtl2syn_verification.rpt
\end{lstlisting}
\end{block}
\end{frame}

\renewcommand{\assubtitle}{Running LEC}
\begin{frame}[fragile]
\frametitle{\astitle}
\begin{block}{\assubtitle}
\begin{lstlisting}[style=sh]
cd RWU_RV/RV_NoPipeline
./scripts/lec/lec.sh                  # rtl2syn, then syn2pnr
./scripts/lec/lec.sh -m rtl2syn       # one check only
./scripts/lec/lec.sh -m rtl2syn -g    # with the Genus dofile
./scripts/lec/lec.sh -m rtl2syn -i    # GUI stays open (debug)
\end{lstlisting}
\begin{lstlisting}[style=out]
=== [lec] rtl2syn
    equivalent NNNNN, non-equivalent 0, abort 0
PASS  LEC rtl2syn
=== [lec] syn2pnr
    equivalent NNNNN, non-equivalent 0, abort 0
PASS  LEC syn2pnr
\end{lstlisting}
  \begin{itemize}
    \item Needs: \cmd{build\_genus/} and for \cmd{syn2pnr}
      \cmd{build\_innovus/cc\_chip\_top\_routed.v}
  \end{itemize}
\end{block}
\end{frame}

\renewcommand{\assubtitle}{Reports in \texttt{build\_lec/}}
\begin{frame}
\frametitle{\astitle}
\begin{block}{\assubtitle}
  \footnotesize
  \begin{tabularx}{\linewidth}{@{}lY@{}}
    \textbf{file} & \textbf{content}\\\hline
    \cmd{lec\_<mode>.log}              & everything, incl. the compare table\\
    \cmd{lec\_<mode>\_design.rpt}      & designs read, black boxes,
                                         undefined cells\\
    \cmd{lec\_<mode>\_unmapped.rpt}    & unmapped key points by class\\
    \cmd{lec\_<mode>\_verification.rpt} & checklist: modelling, mapping,
                                         compare -- all must be ok\\
    \cmd{lec\_<mode>\_noneq.rpt}       & non-equivalent and abort points\\
    \cmd{lec\_<mode>\_statistics.rpt}  & key point counts per type\\
  \end{tabularx}
  \vspace{4pt}
  \begin{itemize}
    \item Read \cmd{\_design.rpt} first: a black box that should not be
      one (missing cell in the Liberty) hides its whole logic
  \end{itemize}
\end{block}
\end{frame}

%%%%%%%%%%%%%%%%%%%%%%%%%%%%%%%%
% Debug
%%%%%%%%%%%%%%%%%%%%%%%%%%%%%%%%
\section{Debugging}
\renewcommand{\astitle}{Debugging}

\renewcommand{\assubtitle}{Non-Equivalent Points}
\begin{frame}[fragile]
\frametitle{\astitle}
\begin{block}{\assubtitle}
  \begin{itemize}
    \item \cmd{./scripts/lec/lec.sh -m rtl2syn -i}: the GUI stays open
  \end{itemize}
\begin{lstlisting}[style=sh]
report_compare_data -class nonequivalent
diagnose -summary                     ;# groups points by cause
diagnose <point>                      ;# counterexample for one point
\end{lstlisting}
  \begin{itemize}
    \item Diagnose shows the input pattern that makes golden and revised
      differ, and the schematic of both cones side by side
    \item Many points with one root cause are typical: fix the first
      group, run again
  \end{itemize}
\end{block}
\end{frame}

\renewcommand{\assubtitle}{Typical Causes}
\begin{frame}
\frametitle{\astitle}
\begin{block}{\assubtitle}
  \footnotesize
  \begin{tabularx}{\linewidth}{@{}lY@{}}
    \textbf{symptom} & \textbf{cause}\\\hline
    all scan FFs non-eq. & SE constraint not applied: pin name in the
                          netlist differs (check \cmd{core/as\_jtag} is
                          not ungrouped)\\
    pads / clock cells   & RTL read without the \cmd{XFAB\_*} defines:
                          behavioural model vs. library cell\\
    some FFs, \cmd{x}   & RTL assigns \cmd{'x} (don't care); Genus picked
                          a value. Avoid \cmd{x} in RTL or set
                          \cmd{set\_x\_conversion}\\
    constant FFs         & \cmd{set\_flatten\_model -seq\_constant}\\
    many not-mapped      & names changed: \cmd{add\_renaming\_rule}, or
                          let Genus write the dofile (\cmd{-g})\\
    aborts               & \cmd{analyze\_abort -compare}; for the
                          multiplier/ALU: \cmd{set\_compare\_effort high}\\
  \end{tabularx}
\end{block}
\end{frame}

\renewcommand{\assubtitle}{The Genus Dofile (\texttt{-g})}
\begin{frame}
\frametitle{\astitle}
\begin{block}{\assubtitle}
  \begin{itemize}
    \item \cmd{01b\_rv64\_syn.tcl} writes it at the end:
      \cmd{write\_do\_lec -golden\_design rtl -revised\_design
      cc\_chip\_top\_syn.v}
    \item Genus knows its own renamings, merged and removed flip-flops
      and writes the matching LEC settings
    \item Synthesis ran before this was added? \cmd{lec.sh -g} writes it
      from \cmd{cc\_chip\_top\_syn.db}
    \item Use: cross-check, or as a template if the own script has many
      not-mapped points. Whether Genus also constrains the internal scan
      enable: VERIFY-LAB
  \end{itemize}
\end{block}
\end{frame}

%%%%%%%%%%%%%%%%%%%%%%%%%%%%%%%%
% Open points
%%%%%%%%%%%%%%%%%%%%%%%%%%%%%%%%
\section{Open Points}
\renewcommand{\astitle}{Open Points}

\renewcommand{\assubtitle}{VERIFY-LAB: Not Yet Run with Conformal}
\begin{frame}
\frametitle{\astitle}
\begin{block}{\assubtitle}
  \begin{itemize}
    \item \cmd{lec -xl -tclmode}: licence option and the Tcl command
      names (\cmd{read\_design}, \cmd{add\_primary\_input -pin -cut})
    \item \cmd{read\_design -define NAME=VALUE} syntax for the cell-name
      defines
    \item Hierarchical pin names \cmd{core/as\_jtag/...} after flattening
    \item Redirection \cmd{> file} in Tcl mode (otherwise everything is in
      \cmd{lec\_<mode>.log})
    \item Layout of the compare table that \cmd{lec.sh} parses
    \item Ends of \cmd{write\_do\_lec} output (\cmd{exit} present?)
  \end{itemize}
\end{block}
\end{frame}

\renewcommand{\assubtitle}{X-FAB NDA}
\begin{frame}
\frametitle{\astitle}
\begin{block}{\assubtitle}
  \begin{itemize}
    \item Netlists, LEC logs, reports and the Genus dofile contain X-FAB
      cell names and paths
    \item \cmd{build\_genus/}, \cmd{build\_innovus/}, \cmd{build\_lec/}
      are in \cmd{.gitignore}: never commit them, never copy them out of
      the lab
    \item Check \cmd{git status} before every commit
  \end{itemize}
\end{block}
\end{frame}

\renewcommand{\assubtitle}{Next Steps}
\begin{frame}
\frametitle{\astitle}
\begin{block}{\assubtitle}
  \begin{itemize}
    \item Property checking without Jasper: SymbiYosys (open source) with
      SVA, first on small blocks (TAP FSM, memory arbiter, FIFO)
    \item riscv-formal: the RV64I core against the ISA (needs the RVFI
      trace port)
    \item Physical verification: DRC, LVS, ERC
  \end{itemize}
\end{block}
\end{frame}

%%%%%%%%%%%%%%%%%%%%%%%%%%%%%%%%
% Q&A
%%%%%%%%%%%%%%%%%%%%%%%%%%%%%%%%
\section{Q\&A}
\renewcommand{\astitle}{Q\&A}
\begin{frame}
\frametitle{\astitle}
\vfill
\centering
\begin{Huge}\textcolor{red}{Questions?}\end{Huge}
\vfill
\end{frame}

\end{document}
